% ============================================================
% Prueba Técnica — Resultados
% Choque petrolero de 2022 y subsidios a combustibles fósiles
% en América Latina y el Caribe
% Candidato: Daniel Mendivelso
% ------------------------------------------------------------
% Compilar:  xelatex -interaction=nonstopmode Prueba_Tecnica_Resultados.tex
% Para ver las NOTAS DEL EXPOSITOR, descomentar la línea:
%     \setbeameroption{show notes}
% (o "show notes on second screen" para presentar con notas).
% ============================================================
\documentclass[aspectratio=169,10pt]{beamer}

% --- Tipografía Unicode (XeLaTeX) ---
% IMPORTANTE: compilar con xelatex. fontspec maneja bien los signos
% ¿ ¡ á é í ... (con inputenc/fontenc el ¿ se renderiza como £).
\usepackage{fontspec}
\setsansfont{texgyreheros}[              % clon de Helvetica, working-paper limpio
  Extension      = .otf,
  UprightFont    = *-regular,
  ItalicFont     = *-italic,
  BoldFont       = *-bold,
  BoldItalicFont = *-bolditalic,
  Scale          = 0.92]
\renewcommand{\familydefault}{\sfdefault}
\usepackage[spanish]{babel}
\usepackage{microtype}

% --- Paquetes ---
\usepackage{amsmath}
\usepackage{graphicx}
\usepackage{booktabs}
\usepackage{xcolor}

% ============================================================
% PALETA "Gris carbón minimal" (sin marca institucional)
% ============================================================
\definecolor{Carbon}{HTML}{333333}     % acento principal
\definecolor{GrisMedio}{HTML}{6E6E6E}   % detalle / secundario
\definecolor{AcentoRojo}{HTML}{B00020}  % UNA cifra clave por slide
\definecolor{GrisClaro}{HTML}{EDEDED}   % fondos sutiles

% --- Tema limpio construido a mano ---
\usetheme{default}
\setbeamertemplate{navigation symbols}{}
\setbeamertemplate{itemize items}[circle]
\setbeamertemplate{itemize subitem}[triangle]

% Colores del tema
\setbeamercolor{structure}{fg=Carbon}
\setbeamercolor{frametitle}{fg=Carbon}
\setbeamercolor{title}{fg=Carbon}
\setbeamercolor{subtitle}{fg=GrisMedio}
\setbeamercolor{author}{fg=GrisMedio}
\setbeamercolor{date}{fg=GrisMedio}
\setbeamercolor{normal text}{fg=Carbon}
\setbeamercolor{itemize item}{fg=Carbon}
\setbeamercolor{itemize subitem}{fg=GrisMedio}
% Énfasis (\alert) en carbón negrita; el rojo se reserva para \key (UNA cifra por slide)
\setbeamercolor{alerted text}{fg=Carbon}
\setbeamerfont{alerted text}{series=\bfseries}
\setbeamercolor{block title}{bg=GrisClaro,fg=Carbon}
\setbeamercolor{block body}{bg=GrisClaro!40,fg=Carbon}

\setbeamerfont{frametitle}{series=\bfseries}
\setbeamerfont{title}{series=\bfseries,size=\LARGE}
\setbeamerfont{block title}{series=\bfseries}

% --- Frametitle: título en carbón + regla fina roja debajo ---
\setbeamertemplate{frametitle}{%
  \vspace{0.6em}%
  \begin{beamercolorbox}[wd=\paperwidth,leftskip=0.6cm,rightskip=0.6cm]{frametitle}%
    {\bfseries\insertframetitle}\par%
    \vspace{0.15em}%
    {\color{AcentoRojo}\rule{2.2cm}{1.2pt}}%
  \end{beamercolorbox}%
  \vspace{0.2em}%
}

% --- Footline minimalista en gris ---
\setbeamertemplate{footline}{%
  \hbox{%
    \begin{beamercolorbox}[wd=\paperwidth,ht=2.2ex,dp=1ex,leftskip=0.6cm,rightskip=0.6cm]{}%
      \usebeamerfont{date}\color{GrisMedio}\scriptsize
      \hfill \insertframenumber%
    \end{beamercolorbox}%
  }%
}

% --- Comando para resaltar la cifra clave (rojo, negrita) ---
\newcommand{\key}[1]{{\bfseries\color{AcentoRojo}#1}}

% --- Atajos de ruta a los assets ---
% figures/cropped/  = figuras 8-bit RGB con la banda de notas recortada
%                     (la nota va como texto LaTeX bajo cada figura).
% figures/compiled/ = figuras 8-bit RGB completas (con su nota quemada),
%                     usadas en los slides de respaldo.
% Las originales de 16-bit no renderizan bajo XeLaTeX/xdvipdfmx.
\newcommand{\fig}[1]{figures/cropped/#1}
\newcommand{\figfull}[1]{figures/compiled/#1}
\newcommand{\tab}[1]{tables/#1}

% Nota al pie de figura: gris, pequeña, JUSTIFICADA A LA IZQUIERDA,
% interlineado compacto. Va dentro de la columna (hereda su ancho).
\newcommand{\fignota}[1]{\par\vspace{0.2em}%
  {\scriptsize\color{GrisMedio}\setlength{\baselineskip}{0.85\baselineskip}%
   \raggedright #1\par}}

% Igual, pero para usar dentro de un \begin{center} (figura full-width):
% envuelve la nota en un minipage del ancho del texto, alineada a la izquierda.
\newcommand{\fignotaC}[1]{\par\vspace{0.2em}%
  \begin{minipage}{0.92\linewidth}%
    {\scriptsize\color{GrisMedio}\setlength{\baselineskip}{0.85\baselineskip}%
     \raggedright #1\par}%
  \end{minipage}}

% Caption tipo paper para figuras/tablas: "Figura N. Título descriptivo".
% Va ENCIMA del \includegraphics. #1 = etiqueta (Figura 1 / Tabla 5), #2 = título breve.
% La numeración sigue el nombre de archivo (fig1..5, tab1,2,4,5,6), no el orden de aparición.
\newcommand{\captitulo}[2]{%
  {\footnotesize\color{Carbon}\textbf{#1.}~#2\par}\vspace{0.2em}}

% --- Divisor de sección (reusa el patrón de "Material de respaldo") ---
% #1 = número de bloque, #2 = título, #3 = subtítulo. Frame [plain] sin footline.
\newcommand{\divisor}[3]{%
  {\setbeamertemplate{footline}{}%
   \begin{frame}[plain]
     \vfill
     \begin{center}
       {\color{AcentoRojo}\rule{2.6cm}{1.2pt}}\\[0.7em]
       {\scriptsize\color{GrisMedio}Bloque #1 \,/\, 7}\\[0.6em]
       {\Large\color{Carbon}\bfseries #2}\\[0.45em]
       {\color{GrisMedio}#3}\\[0.7em]
       {\color{AcentoRojo}\rule{2.6cm}{1.2pt}}
     \end{center}
     \vfill
   \end{frame}}}

% --- Espaciado de listas más aireado ---
\setlength{\parskip}{2pt}

% ============================================================
% Para PRESENTAR CON NOTAS, descomentar una de estas líneas:
% \setbeameroption{show notes}
% \setbeameroption{show notes on second screen=right}
% ============================================================

% --- Metadatos ---
\title{El choque petrolero de 2022 y los subsidios a combustibles fósiles}
\subtitle{Un experimento natural en América Latina y el Caribe}
\author{Daniel Mendivelso}
\date{2026}

\begin{document}

% ============================================================
% PORTADA
% ============================================================
{
\setbeamertemplate{footline}{} % portada sin footline
\begin{frame}[plain]
  \vfill
  \begin{center}
    {\color{AcentoRojo}\rule{3.5cm}{1.4pt}}\\[1.1em]
    {\usebeamerfont{title}\color{Carbon}\bfseries El choque petrolero de 2022\\[0.2em]
     y los subsidios a combustibles fósiles}\\[2.4em]
    {\large\color{Carbon}\textbf{Daniel Mendivelso}}\\[0.4em]
    {\color{GrisMedio}2026}\\[1.1em]
    {\color{AcentoRojo}\rule{3.5cm}{1.4pt}}
  \end{center}
  \vfill
\end{frame}
}

% ############################################################
% BLOQUE 1 — EL PROBLEMA
% ############################################################
\divisor{1}{El choque de 2022 y la pregunta de política}{El problema}

% --- Slide 1 (existente): El choque de 2022 ---
\begin{frame}{El choque de 2022: un experimento natural sobre los subsidios}
  \begin{itemize}
    \item En 2022 el Brent saltó de \key{USD 71 a 101 por barril (+42\%)} en promedio anual:
          el mayor salto de la serie 2015--2023.
    \item Detonante: la \alert{invasión de Rusia a Ucrania} (feb.\ 2022). Un choque
          \textbf{global, inesperado} y ajeno a la política de subsidios de cualquier país.
    \item Es \textbf{plausiblemente exógeno} a la política de subsidios de la región: ahí está
          la oportunidad de identificación.
  \end{itemize}

  \begin{block}{Pregunta y tesis}
    \textbf{¿Cómo respondieron los subsidios fósiles de la región ante el alza del precio del
    Brent en 2022, y qué implicaciones fiscales tiene esa respuesta?}\\[0.3em]
    \textbf{Tesis:} el choque petrolero de 2022 no afectó a todos los países de la región por igual;
    su impacto fiscal, expresado a través del subsidio explícito a los combustibles, dependió de la
    \textbf{posición de exposición petrolera neta} de cada país.
  \end{block}

  \note{%
    El precio del Brent lo fija el mercado mundial; ningún país de la región lo mueve con su
    política de subsidios. Eso es lo que lo hace exógeno y útil como fuente de variación. La
    pregunta no es solo ``¿subieron los subsidios?'' sino ``¿quién y por qué?'', porque ahí
    está la lección de política. Anticipo la tesis: la exposición petrolera neta del país
    determina la dirección del golpe, y esa es la variable sobre la que construyo todo.}
\end{frame}

% --- Slide 1b (NUEVO): experimento natural + describir vs medir + contexto FMI ---
\begin{frame}{Por qué este episodio permite medir un efecto, no solo describirlo}
  \small
  \begin{block}{El choque como experimento natural}
    El precio internacional lo fija el mercado mundial; \textbf{ningún país de la región lo mueve
    con su política de subsidios}. La variación del tratamiento no viene de las decisiones que
    queremos estudiar, sino de afuera.
  \end{block}

  \textbf{Describir $\neq$ medir.} Ver que el subsidio subió en 2022 \emph{describe} una
  evolución. Pero la pregunta es \textbf{causal} (¿cuánto se debe al choque?), y una pregunta
  causal exige un \textbf{contrafactual}: qué habría pasado sin el choque. Para eso necesito un
  grupo de comparación, no solo una serie en el tiempo. El grupo de comparación son los
  \alert{importadores netos}: países expuestos al mismo choque global pero con mecanismo
  opuesto — para ellos el Brent encarece el costo de subsidiar, no genera renta.\\[0.2em]
  \textbf{Tesis:} el impacto fiscal del choque, medido por el subsidio explícito,
  dependió de la \key{posición de exposición petrolera neta} de cada país.

  \begin{block}{Contexto global (FMI)}
    {\footnotesize Según el FMI, los subsidios fósiles globales alcanzaron \textbf{USD 7 billones
    (7.1\% del PIB mundial)} en 2022, y los subsidios \emph{explícitos} más que se duplicaron
    \textbf{desde 2020} (Black, Liu, Parry \& Vernon, 2023).}
  \end{block}

  \note{%
    Este slide fija la lógica de identificación antes de cualquier dato. La clave de un experimento
    natural: la fuente de variación (el precio mundial) es ajena a la política que estudio. Insisto
    en describir vs.\ medir porque es el error típico: una serie que sube no prueba un efecto;
    necesito el
    contrafactual. El dato global del FMI lo atribuyo explícitamente a ellos, como contexto, no como
    hallazgo mío; y ojo: es ``desde 2020'', no ``en 2022''.}
\end{frame}

% ############################################################
% BLOQUE 2 — LOS DATOS
% ############################################################
\divisor{2}{Fuentes y cobertura del panel}{Los datos}

% --- Slide 3 (existente): Datos ---
\begin{frame}{Datos: panel de 34 países, 2015--2023}
  \small
  Construyo un panel anual de los países de América Latina y el Caribe con cobertura completa,
  combinando tres fuentes que cumplen funciones distintas:
  \begin{itemize}
    \item \textbf{Variable dependiente} (lo que quiero explicar): el subsidio explícito a
          combustibles fósiles, medido como \% del PIB. Viene de la \emph{IMF Fossil Fuel
          Subsidies Database}, la fuente obligatoria, que además separa el componente explícito
          del implícito.
    \item \textbf{Tratamiento} (la fuente de variación): el choque del precio del Brent en 2022
          (EIA), que entra al modelo \textbf{interactuado} con la condición de exportador neto.
          El choque por sí solo no identifica nada; lo hace su interacción con el grupo.
    \item \textbf{Variables fiscales} (deuda pública bruta, balance fiscal e ingreso del
          gobierno general, del \textbf{FMI, World Economic Outlook} vía IMF DataMapper). El choque
          también las afecta: son \textbf{variables de resultado}, no covariables exógenas, así que
          incluirlas en el modelo introduciría \textbf{endogeneidad}. Por eso quedan \textbf{fuera
          de la estimación} y se usan solo en la parte de política.
  \end{itemize}

  \begin{block}{Estructura del panel}
    \key{34 países $\times$ 9 años $=$ 306 observaciones}, balanceado, sin \emph{missing values}.
    \textbf{7 exportadores netos} frente a \textbf{27 importadores netos}.\\[0.3em]
    \textbf{Exportadores netos:} Bolivia, Colombia, Ecuador, Guyana, México, Trinidad y Tobago,
    Venezuela.
  \end{block}

  \note{%
    La base del FMI es la fuente obligatoria y distingue explícito de implícito; me centro en el
    explícito. Las fiscales las traje del World Economic Outlook del FMI (vía IMF DataMapper), no
    del Banco Mundial, porque cubre los 34 países sin missing values. Punto metodológico: las fiscales NO
    entran al modelo del efecto --son variables de resultado, afectadas por el propio choque, así que
    introducirían endogeneidad--. Son solo 7 exportadores: grupo chico, con consecuencias en la precisión
    que reconozco abiertamente.}
\end{frame}

% --- Slide 2b (NUEVO): por qué World Economic Outlook sobre Banco Mundial + limitación de tamaño ---
\begin{frame}{Decisiones de fuente y sus límites}
  \small
  \begin{block}{Por qué el World Economic Outlook (FMI) y no el Banco Mundial}
    La parte de política exige cruzar deuda contra subsidio \textbf{país por país}, lo que requiere
    \textbf{cobertura completa} de los 34 países. El \emph{World Economic Outlook} del FMI la tiene;
    el Banco Mundial presenta \alert{amplios \emph{missing values} en el Caribe}, que impiden el
    cruce cuantitativo y la dejarían en uso meramente anecdótico.\\[0.3em]
    {\footnotesize Indicadores WEO (vía IMF DataMapper), todos \% del PIB: deuda pública bruta
    (\texttt{GGXWDG\_NGDP}), balance fiscal (\texttt{GGXCNL\_NGDP}) e ingreso del gobierno general
    (\texttt{GGR\_G01\_GDP\_PT}).}
  \end{block}

  \begin{block}{Por qué la ventana 2015--2023}
    Define una ventana \textbf{simétrica alrededor del choque}: seis años pre-choque (2015--2021)
    para evaluar el supuesto de tendencias paralelas, el año del tratamiento (2022) y el de
    reversión (2023). 2015 es el primer año que extraigo de la base.
  \end{block}

  \textbf{Limitación que reconozco desde el inicio:} hay solo \key{7 exportadores netos}. Eso afecta
  la \textbf{precisión} del estimador (errores estándar mayores, intervalos más anchos), \alert{no su
  insesgadez} (sigue centrado en el efecto verdadero). Lo retomo en los supuestos y la robustez.

  \note{%
    Aquí justifico dos decisiones de datos que un evaluador preguntaría. Primero, por qué el World
    Economic Outlook del FMI (vía IMF DataMapper) y no el Banco Mundial: no es preferencia, es
    cobertura. El Banco Mundial deja muchos missing values en el Caribe, y sin la serie completa no
    puedo cruzar deuda contra subsidio para todos los países, que es justo lo que pide la parte
    fiscal. Segundo, la ventana 2015--2023: es simétrica alrededor del choque (seis años pre para
    examinar tendencias paralelas, el año del tratamiento y el de reversión), y 2015 es el primer
    año que extraigo de la base. Y dejo claro el límite del grupo reducido: afecta la precisión, no
    la insesgadez. Esa distinción reaparece en los supuestos.}
\end{frame}

% --- Slide 2c (NUEVO): composición del panel — Tabla 2 ---
\begin{frame}{Quiénes componen el panel: los 34 países}
  \leavevmode
  {\footnotesize \textbf{Qué mirar:} el \key{grupo de 7 exportadores} (arriba) frente a los 27
  importadores, y la heterogeneidad de niveles dentro de cada grupo (de Venezuela, con subsidio
  altísimo, a países con casi cero). El tratamiento se define por esta clasificación, no por
  ``país productor''.\par}\vspace{0.1em}
  \begin{center}
  \captitulo{Tabla 2}{Clasificación de los 34 países por exposición petrolera neta y subsidio explícito anual}
  \includegraphics[width=0.90\textwidth,height=0.62\textheight,keepaspectratio]{\tab{tab2_paises.pdf}}
  \end{center}

  \note{%
    Muestro la composición real del panel para que el grupo de tratamiento no sea abstracto. Lo que
    quiero que se vea: son solo 7 exportadores, y dentro de cada grupo hay enorme dispersión de
    niveles. Eso conecta con dos cosas dichas antes: la clasificación es por posición neta (Argentina
    y Brasil, aunque producen, son importadores netos de derivados), y el grupo chico de exportadores
    limita la precisión. Es referencia: no la leo entera, la dejo para quien quiera ver país por país.}
\end{frame}

% ############################################################
% BLOQUE 3 — LAS VARIABLES (y la justificación del método)
% ############################################################
\divisor{3}{Qué medimos y por qué}{Las variables y la estrategia de inferencia}

% --- Slide 2 (existente): Dos canales opuestos ---
\begin{frame}{Dos canales opuestos según la exposición petrolera}
  \footnotesize
  El mismo choque de precios opera por \textbf{dos mecanismos distintos} según si el país
  vende o compra petróleo en el mercado internacional.
  \begin{itemize}\setlength{\itemsep}{2pt}
    \item \textbf{Países que venden petróleo} (exportadores netos: Bolivia, Colombia, Ecuador,
          México, Venezuela\ldots): cuando el Brent sube, sus ingresos fiscales crecen.
          Con más recursos, el gobierno puede \textbf{mantener el precio interno bajo} —
          el subsidio explícito se expande y hay renta para financiarlo.
    \item \textbf{Países que compran petróleo} (importadores netos: el resto de LAC):
          cuando el Brent sube, el \textbf{costo de suministro aumenta}. Si el gobierno
          congela el precio al consumidor, el subsidio crece igual, pero \textbf{sin ingresos
          extra que lo respalden}.
  \end{itemize}
  El modelo estima \textbf{cuánto más} creció el subsidio en los países vendedores respecto
  a los compradores: ese diferencial es el efecto del choque.\vspace{0.3em}

  \textbf{¿Por qué ``neto'' y no ``productor''?} Un país puede extraer crudo y aun así importar
  más derivados de los que exporta — como Argentina o Brasil. Lo que determina si el gobierno
  gana renta con el alza del Brent es la \alert{balanza comercial petrolera neta}, no la
  producción.

  \begin{block}{Por qué el subsidio explícito (y no el implícito)}
    El \key{subsidio explícito} es la brecha entre el precio al consumidor y el costo de
    suministro: reacciona directamente al precio del Brent. El \textbf{implícito} (externalidades
    ambientales e impuestos no cobrados) está dominado por su componente estructural, de ajuste
    lento, y en la práctica apenas reaccionó al choque. Por eso la variable de resultado es el
    explícito.
  \end{block}

  \note{%
    Aquí justifico la variable de tratamiento. Clasifico por exportador/importador \emph{neto}
    y no por ``país productor'' porque lo que importa para el efecto fiscal es la posición neta:
    un país puede producir petróleo y aun así ser importador neto de derivados. El explícito es
    precio-sensible (una brecha de precios); el implícito está dominado por su componente estructural.
    Si el efecto aparece en el explícito y es nulo en el implícito, descarta un artefacto contable y
    refuerza el canal del precio: por eso lo uso como placebo más adelante. Caveat honesto: el implícito
    incluye IVA no cobrado, que sí escala con el precio, así que es un placebo imperfecto; lo leo como
    evidencia de apoyo, no como prueba.}
\end{frame}

% --- Slide 3c (NUEVO): por qué DiD, como cadena lógica ---
\begin{frame}{Por qué diferencias en diferencias}
  \small
  La elección del método se deriva del problema y de la estructura de los datos. La cadena lógica:
  \begin{enumerate}
    \item La pregunta es \textbf{causal}: exige estimar un \textbf{contrafactual}, la trayectoria
          que habrían seguido los expuestos sin el choque.
    \item Comparar a los expuestos \textbf{antes y después} no aísla el choque: entre 2021 y 2022
          también cambiaron otros factores globales (rebote pos-pandemia, inflación), y esa
          comparación se los atribuiría todos al choque.
    \item Comparar exportadores con importadores \textbf{solo en 2022} tampoco lo aísla: los grupos
          ya diferían antes del choque (en 2021 los exportadores subsidiaban más), y esa diferencia
          previa se confundiría con el efecto.
    \item Los datos tienen \key{dos dimensiones, tiempo y grupo}. Los importadores atravesaron el
          mismo período sin estar expuestos al canal de la renta petrolera.
    \item Tomarlos como contrafactual y restar a los expuestos lo común al período
          \alert{es, precisamente, diferencias en diferencias}.
  \end{enumerate}

  \note{%
    Este es el slide bisagra entre describir y modelar. La idea que quiero transmitir: el DiD se
    deriva de la pregunta causal y de la estructura de los datos, no se impone. Una comparación
    antes/después confunde el efecto con todo lo que cambió en 2022; un corte transversal lo confunde
    con los niveles que los grupos ya tenían antes. Hacen falta las dos diferencias a la vez. El grupo
    de control surge de los propios datos: los importadores atravesaron el mismo período sin el canal
    de la renta. Restar a los expuestos lo común al período es, literalmente, la doble diferencia.}
\end{frame}

% ############################################################
% BLOQUE 4 — ESTADÍSTICA DESCRIPTIVA
% ############################################################
\divisor{4}{Qué muestran los datos antes del modelo}{Estadística descriptiva}

% --- Slide 6 (existente): Ruptura 2022 — Figura 1 ---
\begin{frame}{Subsidio explícito en 2022: incremento concentrado en exportadores netos}
  \begin{columns}[T]
    \begin{column}{0.46\textwidth}
      \small
      La figura traza el subsidio explícito en \% del PIB, año a año, por grupo.\\[0.25em]
      \textbf{El incremento se concentra en los exportadores:} de
      \textbf{1.40\% a 2.54\% del PIB} (2021$\to$2022), \key{+1.14 pp}.
      Los importadores apenas se mueven: de \textbf{0.55\% a 0.66\%} (+0.12 pp).\\[0.25em]
      \textbf{En el agregado} el subsidio pasa de 0.85\% a 1.33\% del PIB, pero ese
      promedio mezcla los dos grupos y \alert{oculta que el movimiento proviene
      casi exclusivamente de los exportadores}.\\[0.25em]
      Esa asimetría es la que el modelo formaliza y cuantifica.
    \end{column}
    \begin{column}{0.54\textwidth}
      \centering
      \captitulo{Figura 1}{Subsidio explícito por grupo y año, 2015--2023}
      \includegraphics[width=\linewidth,height=0.60\textheight,keepaspectratio]{\fig{fig1_ruptura.png}}
      \fignota{\textit{Cómo leerla:} cada panel es la suma anual del subsidio explícito;
      (a)--(b) en USD miles de millones, (c)--(d) en \% del PIB. Fuente: IMF FFS Database.}
    \end{column}
  \end{columns}

  \note{%
    Empiezo por la evidencia descriptiva antes de cualquier modelo. El mensaje no es que el subsidio
    subió en 2022 (el agregado lo hace, de 0.85 a 1.33\% del PIB), sino que ese salto se concentra en
    los exportadores (1.40 a 2.54) y casi no toca a los importadores (0.55 a 0.66). Por eso pongo el
    contraste por grupo primero y el agregado como contexto: el promedio mezcla ambos y oculta el
    patrón. Los montos en dólares y el componente implícito (placebo) van en la Tabla 1, el slide
    siguiente. Aquí agrego ponderando por PIB; el modelo trata cada país como una unidad, por eso los
    niveles del DiD difieren, aunque el patrón es el mismo. La descriptiva es el diagnóstico que
    motivó el método.}
\end{frame}

% --- Slide 6b (NUEVO): la ruptura en números — Tabla 1 ---
\begin{frame}{Subsidio explícito e implícito: respuestas asimétricas al choque}
  \leavevmode
  {\footnotesize \textbf{Qué mirar:} la fila \emph{Explícito (\% PIB)}, columna \textbf{2022}
  frente a 2021, en cada panel. El total salta de 0.85 a \key{1.33}; el explícito de exportadores
  de 1.40 a \textbf{2.54} y el de importadores apenas de 0.55 a \textbf{0.66}. El \emph{implícito}
  varía en niveles, pero el \textbf{efecto del choque sobre él es nulo} ($-0.32$, no significativo):
  por eso sirve de placebo.\par}\vspace{0.1em}
  \begin{center}
  \captitulo{Tabla 1}{Subsidio explícito, implícito y total por grupo (USD y \% del PIB), 2015--2023}
  \includegraphics[width=0.92\textwidth,height=0.60\textheight,keepaspectratio]{\tab{tab1_descriptiva.pdf}}
  \end{center}

  \note{%
    Este slide da los números exactos detrás de la figura de ruptura, para quien quiera leerlos. La
    lectura guiada: fijarse en la fila del explícito en \% del PIB, columna 2022, en los tres
    paneles. El salto del total (0.85 a 1.33) y, sobre todo, el contraste exportadores (1.40 a 2.54)
    vs.\ importadores (0.55 a 0.66). Y noto que el efecto estimado sobre el implícito es nulo
    ($-0.32$, no significativo): aunque su serie varía en niveles, no responde al choque en el
    modelo, lo que anticipa su uso como placebo. La tabla agrega ponderando por PIB; el modelo
    trata cada país como unidad.}
\end{frame}

% --- Slide 7 (existente): Co-movimiento Brent — Figura 2 ---
\begin{frame}{Vínculo precio–subsidio: evidencia descriptiva y sus límites}
  \begin{itemize}\small
    \item \textbf{Mecanismo:} cuando el precio internacional sube y \textbf{no se traslada}
          al consumidor, la brecha entre costo de suministro y precio regulado se amplía —
          es decir, el subsidio explícito crece de forma casi mecánica.
    \item \textbf{Evidencia:} Brent y subsidio explícito agregado registran una
          \key{correlación de 0.75} en 2015--2023, con un pico conjunto en 2022.
    \item \textbf{Límite:} la correlación es agregada — combina exportadores (donde el alza
          del Brent infla la renta que financia el subsidio) e importadores (donde encarece el
          costo). \alert{Describe el vínculo; no identifica el efecto causal.}
  \end{itemize}
  \vspace{-0.6em}
  \begin{center}
    \captitulo{Figura 2}{Co-movimiento del Brent y el subsidio explícito agregado (doble eje), 2015--2023}
    \includegraphics[height=0.35\textheight]{\fig{fig2_brent.png}}
    \fignotaC{\textit{Cómo leerla:} doble eje, Brent (azul, izq.) y subsidio explícito de
    LAC (naranja, der.), 2015--2023. Importa el co-movimiento, no el nivel.
    Fuente: IMF FFS; EIA.}
  \end{center}

  \note{%
    Esta figura conecta tratamiento (Brent) y resultado (subsidio) de forma visual y honesta. La
    correlación de 0.75 es fuerte, pero correlación agregada $\neq$ causa: en el mismo número están
    sumados dos grupos con mecanismos opuestos. La uso para motivar, no para concluir. Marca la
    diferencia entre describir e identificar, que es lo que hace la parte siguiente. Las
    descriptivas son el diagnóstico que dictó el método.}
\end{frame}

% --- Slide 8b (NUEVO): heterogeneidad país a país — Figura 3 ---
\begin{frame}{Cambio del subsidio explícito por país: distribución asimétrica en 2022}
  \begin{columns}[T]
    \begin{column}{0.50\textwidth}
      \small
      Cambio del subsidio explícito por país (pp del PIB), 2022 vs.\ promedio 2015--2019.
      \begin{itemize}
        \item \textbf{Mayores incrementos:} \key{Bolivia (+5.47), Venezuela (+5.25),
              Colombia (+2.45)}. Los exportadores netos (naranja) dominan el extremo derecho.
        \item \textbf{Variación intra-grupo:} Argentina, importador neto, sube +2.34 pp —
              por encima de varios exportadores. El DiD mide el diferencial \emph{promedio}
              entre grupos, no la respuesta de cada país.
        \item \textbf{Contraejemplo:} Guyana, exportador neto, \alert{redujo} su subsidio
              ($-0.23$ pp): el patrón de grupo no opera de forma uniforme.
      \end{itemize}
    \end{column}
    \begin{column}{0.50\textwidth}
      \centering
      \captitulo{Figura 3}{Cambio del subsidio por país, 2022 vs.\ 2015--2019}
      \includegraphics[width=\linewidth,height=0.58\textheight,keepaspectratio]{\fig{fig3_impacto.png}}
      \fignota{\textit{Cómo leerla:} cambio del subsidio explícito por país (pp del PIB),
      2022 vs.\ 2015--2019; naranja $=$ exportador, azul $=$ importador. Fuente: IMF FFS.}
    \end{column}
  \end{columns}

  \note{%
    Cierro la descriptiva mostrando la distribución país a país. Tres lecturas: los exportadores
    dominan el extremo derecho (Bolivia, Venezuela, Colombia), pero Argentina —importador— también
    sube fuerte (+2.34 pp), lo que muestra que el subsidio es precio-sensible en ambos grupos; el
    DiD captura el diferencial promedio, no la respuesta individual. Guyana es el contraejemplo
    clave: exportador que reduce su subsidio, consistente con que recibió ingresos petroleros
    extraordinarios y los destinó a otro fin. Esta heterogeneidad justifica dos decisiones: en el
    modelo trato cada país como unidad (no pondero por PIB), y en política cruzo el efecto estimado
    con el espacio fiscal de cada economía en lugar de una recomendación uniforme (Bloque 7).}
\end{frame}

% ############################################################
% BLOQUE 5 — EL MODELO
% ############################################################
\divisor{5}{Diferencias en diferencias: especificación y efecto}{El modelo}

% --- Slide 4 (existente): Estrategia DiD ---
\begin{frame}{Diferencias en diferencias: especificación TWFE}

  La estrategia compara el \textbf{cambio} del subsidio en exportadores netos
  contra el \textbf{cambio} en importadores netos, del período pre-choque
  (2015--2021) al año del choque (2022).
  \begin{block}{Especificación TWFE}
    \[
      \text{Subsidio}_{it} = \textcolor{AcentoRojo}{\boldsymbol{\beta_3}}\,
      (\text{Post2022}_t \times \text{Exportador}_i)
      + \alpha_i + \lambda_t + \varepsilon_{it}
    \]
    \begin{itemize}\small
      \item $\alpha_i$: absorben lo invariante del país, incluido ser exportador.
      \item $\lambda_t$: absorben lo común a cada año, incluido el nivel del Brent.
      \item $\beta_3$: cambio diferencial del subsidio de los exportadores atribuible al choque.
    \end{itemize}
  \end{block}

  \begin{itemize}\small
    \item \textbf{Por qué la interacción:} $\text{Post2022}_t$ fecha el choque
          y $\text{Exportador}_i$ identifica la posición petrolera neta. Su producto
          recoge la respuesta del subsidio al choque entre las economías que captan
          renta petrolera. La identificación proviene \textbf{exclusivamente} de
          esta interacción: el nivel del Brent lo absorbe $\lambda_t$.
    \item \textbf{Referencia:} DiD 2$\times$2 canónico (Angrist \& Pischke, 2009).
          La validez descansa en tendencias paralelas pre-2022, verificadas con
          el placebo (Bloque 6).
  \end{itemize}

  \note{%
    El truco del DiD: no basta ver que el subsidio de los exportadores subió --pudo subir por mil
    razones--. Lo comparo contra un control (importadores) que vivió el mismo año sin el canal de
    la renta petrolera. La doble diferencia limpia los factores comunes. Punto central: el Brent
    es el mismo para todos en cada año $\Rightarrow$ varianza cero dentro de año $\Rightarrow$ lo
    absorbe $\lambda_t$ y no lo identifico solo. La identificación viene EXCLUSIVAMENTE de la
    interacción. Si preguntan qué representa la interacción: Post2022 fecha el choque de precios y
    Exportador identifica la posición petrolera neta, que es la que expone al país a la renta. El
    producto mide entonces cómo responde el subsidio al choque entre las economías que captan esa
    renta. La clasificación es por posición neta, de modo que un país puede producir crudo y aun así
    quedar en el grupo de control si importa más derivados de los que exporta. Subrayo el
    la especificación se fijó antes de estimar.}
\end{frame}

% --- Slide 9 (existente): Efecto central — Tabla 4 ---
\begin{frame}{Efecto del choque sobre el subsidio explícito: +1.80 pp del PIB}
  \begin{columns}[T]
    \begin{column}{0.50\textwidth}
      \footnotesize
      El DiD 2$\times$2 (col.\ 1) se construye de cuatro medias grupales:
      \begin{itemize}
        \item Exportadores: $3.30 \to 5.56$\,\% del PIB\ \ ($+2.26$ pp)
        \item Importadores: $0.43 \to 0.88$\,\% del PIB\ \ ($+0.45$ pp)
        \item Doble diferencia $=$ \key{+1.80 pp del PIB}
      \end{itemize}
      \vspace{0.2em}
      La col.\ (2) replica ese resultado con efectos fijos: $\hat{\beta}_3 = +1.79$ pp
      (EE 0.92), significativo a $\alpha = 0.10$ ($p \approx 0.06$). La magnitud es
      sustantiva; la precisión es limitada por el reducido número de exportadores ($n = 7$).
      \vspace{0.2em}

      \textbf{Placebo (cols.\ 3--4):} el efecto aparece en el subsidio explícito
      ($+1.79$\,pp, $p \approx 0.06$) pero es nulo en el implícito ($-0.32$\,pp,
      \alert{no significativo}). Eso es lo esperado si el canal es el precio: el
      explícito es sensible al Brent; el implícito está dominado por factores
      estructurales de ajuste lento.
    \end{column}
    \begin{column}{0.50\textwidth}
      \centering
      \captitulo{Tabla 4}{Efecto del choque (DiD/TWFE): explícito, placebo y total}
      \includegraphics[width=\linewidth,height=0.52\textheight,keepaspectratio]{\tab{tab4_modelo.pdf}}
      \fignota{\textit{Cómo leerla:} cada columna es una especificación; la fila superior
      es $\hat{\beta}_3$ (pp del PIB) con su EE agrupado por país entre paréntesis.
      Fuente: IMF FFS; EIA.}
    \end{column}
  \end{columns}

  \note{%
    Muestro primero el DiD calculado directamente --las cuatro medias-- para que se vea el origen
    del número antes de la especificación con efectos fijos: exportadores +2.25, importadores +0.45,
    diferencia 1.80. El modelo con EF da prácticamente lo mismo (1.79): el resultado no depende de la
    especificación. Sobre la significancia, soy explícito: p $\approx$ 0.06, marginal al 10\%; con 7
    exportadores la precisión es limitada. El argumento decisivo es el placebo: efecto en el
    explícito, nulo en el implícito. Si fuera un artefacto aparecería en ambos componentes.}
\end{frame}

% ############################################################
% BLOQUE 6 — PRUEBAS DE LOS SUPUESTOS
% ############################################################
\divisor{6}{¿Se sostiene la identificación?}{Pruebas de los supuestos}

% --- Slide 5a (existente): Supuestos (1) tendencias paralelas ---
\begin{frame}{Supuestos de identificación (1): tendencias paralelas}
  \footnotesize
  \textbf{Tendencias paralelas:} \emph{sin el choque, ambos grupos habrían evolucionado igual.}
  \[
    \mathbb{E}[Y^0_{i,2022}-Y^0_{i,2021}\mid \text{Exp}{=}1]
    = \mathbb{E}[Y^0_{i,2022}-Y^0_{i,2021}\mid \text{Exp}{=}0]
  \]
  \begin{block}{El supuesto se sostiene}
    \begin{itemize}
      \item Los coeficientes pre-choque (2015--2020) alternan de signo
            ($+0.04,\,-0.41,\,-0.09,\,+0.95,\,+0.26,\,-0.77$) y sus \textbf{IC 95\% incluyen el cero}:
            no se observa una tendencia divergente que escale hacia 2022.
      \item El test conjunto de Wald \alert{rechaza} la nulidad estricta
            (\key{F = 4.09, p < 0.001}), pero el rechazo proviene de la \textbf{volatilidad}
            del grupo (7 exportadores expuestos a los shocks de 2018 y COVID-2020),
            \textbf{no de una tendencia sistemática} previa al choque.
    \end{itemize}
  \end{block}
  \centering\footnotesize\color{GrisMedio}
  La volatilidad reduce la precisión de las estimaciones; una tendencia divergente introduciría sesgo. Lo primero se observa; lo segundo, no.

  \note{%
    Esto es un examen, no una declaración de victoria. Lo que importa para que el DiD no esté
    sesgado es que no haya una \emph{tendencia} divergente previa, y no la hay: los coeficientes
    alternan de signo y sus intervalos cruzan el cero. El Wald rechaza que sean exactamente cero
    (p < 0.001), pero distingo dos cosas: el rechazo es por volatilidad --vaivenes de un grupo
    chico golpeado por 2018 y la pandemia--, no por una pendiente hacia 2022. La volatilidad me
    cuesta precisión, no insesgadez.}
\end{frame}

% --- Slide 8 (existente): Event study — Figura 4 (evidencia del supuesto 1) ---
\begin{frame}{Event study: ausencia de tendencia divergente previa al choque}
  \begin{itemize}\small
    \item Cada punto es el efecto diferencial exportadores--importadores en un año,
          con 2021 como referencia. Los coeficientes pre-2022 oscilan en torno al cero
          sin patrón acumulativo: los grupos no venían separándose antes del choque.
    \item El quiebre se concentra en \key{2022 ($\hat{\beta} = +1.79$ pp)}: es una
          ruptura puntual, no la prolongación de una tendencia preexistente.
          En 2023 el coeficiente cae a $+0.14$ pp (no significativo), consistente
          con la reversión del precio del Brent.
    \item El test conjunto de Wald rechaza la nulidad estricta (F = 4.09, p < 0.001)
          por la \textbf{volatilidad} del grupo de exportadores ($n = 7$),
          \textbf{no por una pendiente sistemática} hacia 2022.
  \end{itemize}
  \vspace{-0.6em}
  \begin{center}
    \captitulo{Figura 4}{Event study: efecto diferencial año a año (base 2021)}
    \includegraphics[width=0.82\textwidth,height=0.36\textheight,keepaspectratio]{\fig{fig4_eventstudy.png}}
    \fignotaC{\textit{Cómo leerla:} efecto diferencial (pp del PIB), año a año, base 2021;
    barras $=$ IC 95\%. Fuente: IMF FFS Database.}
  \end{center}

  \note{%
    Aquí está la honestidad metodológica. Distinción crucial: una tendencia sistemática pre-choque
    (coeficientes que escalan) sesgaría el estimador; volatilidad alrededor de cero solo infla la
    incertidumbre. El test rechaza por lo segundo, no por lo primero: lo demuestro mostrando que
    los signos alternan. Argumento decisivo: la brecha no venía creciendo, aparece en 2022 y
    desaparece en 2023 cuando el Brent baja. Por eso 2023 sale del estimador estático (es reversión)
    pero lo conservo en el event study para mostrar que el efecto se revierte.}
\end{frame}

% --- Slide 5b (existente): Supuestos (2) no anticipación ---
\begin{frame}{Supuestos de identificación (2): no anticipación}
  \textbf{No anticipación:} \emph{el subsidio de 2021 no reacciona a un choque de 2022 aún inexistente.}

  \begin{block}{Argumento: el supuesto descansa en la exogeneidad del choque}
    \begin{itemize}
      \item El supuesto no admite un test estándar en un diseño 2$\times$2.
      \item Argumento institucional: \textbf{la invasión de Ucrania (febrero de 2022) fue
            imprevista}; no existe evidencia de que algún gobierno reposicionara sus
            subsidios en 2021 anticipando el alza del precio del petróleo.
      \item La plausibilidad descansa en la \textbf{exogeneidad del choque} respecto
            a la política de subsidios de cualquier país de LAC de forma individual.
    \end{itemize}
  \end{block}

  \note{%
    Este supuesto no se testea en un 2$\times$2, así que no finjo que lo validé con datos. Lo
    defiendo con el diseño: la guerra fue un shock inesperado; es implausible que alguien ajustara
    subsidios en 2021 por algo que no sabía que iba a pasar. Esa es la fortaleza de usar un choque
    exógeno. Cierro recordando que el Brent solo entra vía la interacción, porque su nivel lo come
    el efecto fijo de año.}
\end{frame}

% --- Slide 10 (existente): Robustez (leave-one-out) — Tabla 6 ---
\begin{frame}{Análisis de sensibilidad: $\hat{\beta}_3$ ante la exclusión de casos extremos}
  \begin{columns}[T]
    \begin{column}{0.52\textwidth}
      \small
      Venezuela (subsidio de 18.1\,\% del PIB) y Surinam (8.5\,\% del PIB en el grupo
      de control) son valores atípicos que podrían sesgar el resultado. La tabla
      re-estima $\hat{\beta}_3$ excluyéndolos:
      \begin{itemize}\footnotesize
        \item Muestra completa: $1.79\dagger$
        \item Sin Venezuela: $1.15$ (sin significancia a $\alpha=0.10$; el efecto
              se reduce porque Venezuela \alert{amplifica} pero no genera el resultado).
        \item Sin Surinam: $1.98^{*}$ (sube al retirar un importador con subsidio
              alto que introduce heterogeneidad en el grupo de control).
        \item Sin ambos: $1.33\dagger$
      \end{itemize}
      \vspace{0.2em}
      El \textit{leave-one-out} sobre los 7 exportadores mantiene $\hat{\beta}_3$
      positivo en todo el rango \key{[1.15, 2.21]}: el efecto no depende de un
      único país.
    \end{column}
    \begin{column}{0.48\textwidth}
      \centering
      \captitulo{Tabla 6}{Robustez: $\beta_3$ excluyendo casos extremos}
      \includegraphics[width=\linewidth,height=0.48\textheight,keepaspectratio]{\tab{tab6_robustez.pdf}}
      \fignota{\textit{Cómo leerla:} cada columna es la misma especificación TWFE sobre una
      submuestra distinta. Si $\beta_3$ no cambia de signo, el efecto no depende de ese país.}
    \end{column}
  \end{columns}

  \note{%
    Anticipo la crítica: con un valor extremo como Venezuela, cabe sospechar que el resultado se
    debe solo a ese país. En lugar de excluirla de forma discrecional --una selección sobre el
    resultado, que sería un error-- aplico el leave-one-out a los siete. El efecto nunca cambia de
    signo y queda entre 1.15 y 2.21. Venezuela lo amplifica pero no lo genera. Un resultado relevante:
    excluir a Surinam \emph{eleva} el efecto, porque es un importador con subsidio muy alto que
    introduce heterogeneidad en el grupo de control. Mantengo la muestra completa como especificación
    principal y reporto la sensibilidad por separado.}
\end{frame}

% ############################################################
% BLOQUE 7 — RECOMENDACIONES DE POLÍTICA PÚBLICA
% ############################################################
\divisor{7}{Del coeficiente a la recomendación fiscal}{Recomendaciones de política pública}

% --- Slide 11 (existente): Matriz de política fiscal — Figura 5 ---
\begin{frame}{Matriz de política fiscal: subsidio y espacio fiscal por país (2022)}
  \begin{columns}[T]
    \begin{column}{0.42\textwidth}
      \footnotesize
      El eje horizontal mide el subsidio explícito (\% del PIB); el eje vertical,
      la deuda pública bruta (\% del PIB). Las medianas de la muestra
      (subsidio: 0.81\%, deuda: 64.1\%) dividen el plano en cuatro cuadrantes.\\[0.4em]
      \textbf{Cuadrante superior derecho} (subsidio alto y deuda alta):
      \begin{itemize}\scriptsize
        \item \key{Reforma urgente:} Venezuela (subsidio 18.1\,\%, deuda 164\,\%,
              balance $-5.3$\,\%), Surinam (8.5\,\%, 112\,\%, $-2.6$\,\%),
              Bolivia (8.0\,\%, 69\,\%, $-6.1$\,\%), Argentina.
        \item \textbf{Reforma gradual} (cuadrante inferior derecho): Ecuador,
              Trinidad y Tobago — subsidio alto pero con mayor margen fiscal.
      \end{itemize}
      \vspace{0.2em}
      El tamaño del punto refleja el costo absoluto: aprox.\ 26 mil millones USD
      en México, aprox.\ 6.8 mil millones en Colombia.
    \end{column}
    \begin{column}{0.58\textwidth}
      \centering
      \captitulo{Figura 5}{Matriz subsidio--deuda y cuadrantes de política (2022)}
      \includegraphics[width=\linewidth,height=0.55\textheight,keepaspectratio]{\fig{fig5_matriz_fiscal.png}}
      \fignota{\textit{Cómo leerla:} subsidio (eje X) vs.\ deuda (eje Y), \% del PIB 2022;
      tamaño $=$ costo contrafactual del choque (efecto medio de 1.79\% del PIB aplicado por país,
      no observado). Arriba-derecha $=$ máxima presión fiscal.
      Fuente: IMF Fossil Fuel Subsidies Database; FMI, World Economic Outlook (vía IMF DataMapper).}
    \end{column}
  \end{columns}

  \note{%
    Esta es la sección de política, lo que el ejercicio pide al final. ¿Por qué deuda y no balance?
    Porque el espacio fiscal en la literatura FMI/BID se ancla en el stock de deuda, no en el déficit
    de un año. El cuadrante crítico lo componen Venezuela, Surinam, Bolivia y Argentina. Dos
    aclaraciones: el costo en dólares es el efecto \emph{medio} aplicado a cada economía (un
    contrafactual, no el costo observado), y para Guyana, que redujo su subsidio, no reporto ese costo porque
    contradiría su dato. La presión fiscal NO es exclusiva de exportadores: Surinam y Argentina,
    importadores, están entre los más expuestos.}
\end{frame}

% --- Slide 11a (NUEVO): cómo se construye la matriz / los cuadrantes ---
\begin{frame}{Construcción de la matriz: ejes, cortes y cuadrantes}
  \small
  La matriz es una herramienta de \textbf{priorización descriptiva}, no un modelo causal. Cubre los
  \textbf{26 países con subsidio explícito $>$ 0.05\,\% del PIB} en 2022. Se construye en tres pasos:
  \begin{enumerate}
    \item \textbf{Dos ejes:} subsidio explícito (2022, \% PIB, horizontal) y deuda pública
          bruta (2022, \% PIB, vertical).
    \item \textbf{Dos cortes:} las \key{medianas de la muestra} (subsidio $= 0.81$\,\%,
          deuda $= 64.1$\,\%) dividen el plano. Se emplean medianas —no medias— por su
          robustez ante valores extremos como Venezuela.
    \item \textbf{Cuatro cuadrantes} según la posición relativa a cada mediana:
  \end{enumerate}
  \vspace{-0.2em}
  \begin{columns}[T]
    \begin{column}{0.49\textwidth}
      \begin{block}{Subsidio alto}
        \footnotesize
        \textbf{+ deuda alta $\Rightarrow$} \alert{Reforma urgente}\\
        \emph{(costo fiscal elevado y margen reducido)}\\[0.3em]
        \textbf{+ deuda baja $\Rightarrow$} \alert{Reforma gradual}\\
        \emph{(costo elevado, pero con margen fiscal)}
      \end{block}
    \end{column}
    \begin{column}{0.49\textwidth}
      \begin{block}{Subsidio bajo}
        \footnotesize
        \textbf{+ deuda alta $\Rightarrow$} \alert{Vigilar}\\
        \emph{(gasto moderado, margen fiscal estrecho)}\\[0.3em]
        \textbf{+ deuda baja $\Rightarrow$} \alert{Sin presión inmediata}\\
        \emph{(ni costo elevado ni restricción fiscal)}
      \end{block}
    \end{column}
  \end{columns}
  \vspace{0.1em}
  {\footnotesize \textbf{Nota metodológica:} el espacio fiscal se aproxima con el
  \textbf{stock de deuda acumulada}, no con el déficit anual, que es volátil e influenciado
  por el ciclo.}

  \note{%
    Antes de mostrar la tabla, explico cómo se arma la matriz para que el cuadrante no parezca
    arbitrario. Tres pasos: dos ejes (subsidio y deuda, ambos \% PIB 2022), dos cortes en las
    medianas (0.81 y 64.1; uso medianas porque son robustas a outliers como Venezuela), y los cuatro
    cuadrantes según el lado de cada mediana. Insisto en que es priorización descriptiva, no causal: te
    dice dónde mirar primero, no qué política exacta aplicar. Y justifico el eje de deuda: el espacio
    fiscal es un tema de stock acumulado, no del déficit de un año. Esto prepara la lectura de la
    Tabla 5.}
\end{frame}

% --- Slide 11b (NUEVO): la matriz en detalle — Tabla 5 ---
\begin{frame}{La matriz, país por país: subsidio, deuda y cuadrante}
  \leavevmode
  {\footnotesize La columna \emph{Cuadrante} resume la priorización. El bloque de
  \key{``Reforma urgente''} (subsidio alto $+$ deuda alta) incluye exportadores (Venezuela,
  Bolivia) e importadores (Surinam, Argentina): la presión fiscal no se organiza por
  clasificación petrolera. La columna \emph{Costo del choque} registra n.a.\ para Guyana,
  cuyo subsidio se redujo en 2022 y para el cual el efecto medio estimado no aplica.\par}\vspace{0.1em}
  \begin{center}
  \captitulo{Tabla 5}{Espacio fiscal por país: subsidio, deuda, balance y cuadrante asignado (2022)}
  \includegraphics[width=0.95\textwidth,height=0.60\textheight,keepaspectratio]{\tab{tab5_fiscal.pdf}}
  \end{center}

  \note{%
    Esta tabla convierte la matriz visual en algo accionable: cada país con su subsidio, su deuda,
    su balance y el cuadrante que le toca. La lectura: la columna de cuadrante es la recomendación.
    Lo importante que la figura no muestra tan claro es el balance fiscal (varios de ``reforma
    urgente'' ya están en déficit fuerte: Bolivia $-6.1$, Venezuela $-5.3$). Y refuerzo el mensaje:
    en ``reforma urgente'' hay exportadores e importadores mezclados, así que la recomendación no se
    organiza por grupo petrolero sino por espacio fiscal. El costo en dólares es contrafactual
    (efecto medio aplicado), por eso es n.a.\ en Guyana, cuyo subsidio se movió en sentido contrario.}
\end{frame}

% --- Slide 7b (NUEVO): regresividad / por qué focalizar ---
\begin{frame}{Por qué la reforma debe compensar de forma focalizada}
  \small
  La reducción de un subsidio eleva el precio del combustible y afecta el ingreso real de los
  hogares. La pregunta de política no es \emph{si} reformar, sino \textbf{cómo proteger a los
  hogares de menor ingreso durante la transición}.

  \begin{block}{El subsidio universal a los combustibles es regresivo (FMI)}
    El subsidio beneficia en proporción al consumo, que se concentra en los hogares de mayor
    ingreso: el \key{quintil más rico capta más del 40\%} de los beneficios, aproximadamente
    6 veces la porción del quintil más pobre (Coady, Flamini \& Sears, 2015).
  \end{block}

  \textbf{Implicación de diseño:} un subsidio universal al precio es un instrumento mal
  focalizado. Sustituirlo por \alert{transferencias focalizadas} a los hogares vulnerables
  \textbf{mejora la equidad} y libera espacio fiscal al eliminar la transferencia implícita
  hacia los quintiles de mayor ingreso. Una reforma bien diseñada no es regresiva.

  \note{%
    Este slide sostiene la recomendación de focalizar. El dato de regresividad lo atribuyo
    explícitamente al FMI (Coady et al.\ 2015), no como hallazgo propio: el 20\% más rico capta más
    del 40\% de los beneficios, unas 6 veces la porción del 20\% más pobre. La lección de diseño: el
    subsidio universal al precio es un instrumento mal focalizado; la transferencia dirigida protege
    a los hogares vulnerables y reduce el costo fiscal a la vez. Eso responde a la objeción de que
    ``reformar perjudica a los más pobres'': los perjudica si se hace mal; bien diseñada, es progresiva.}
\end{frame}

% --- Slide 11c (NUEVO): hoja de ruta por cuadrante ---
\begin{frame}{Recomendaciones diferenciadas por cuadrante fiscal}
  \small
  La recomendación \textbf{se diferencia por cuadrante}, no es uniforme. Se calibra según
  el margen fiscal de cada economía:
  \begin{block}{Acción por cuadrante}
    \footnotesize
    \begin{itemize}
      \item \alert{Reforma urgente} (Venezuela, Surinam, Bolivia, Argentina): ajuste
            \textbf{prioritario} del precio interno con \textbf{compensación focalizada inmediata}.
            El balance fiscal ya es deficitario (\key{Bolivia $-6.1$\,\%, Venezuela $-5.3$\,\%
            del PIB}); el margen para postergar el ajuste es reducido.
      \item \alert{Reforma gradual} (Ecuador, Trinidad y Tobago): la deuda pública es
            relativamente baja, lo que permite una \textbf{secuencia escalonada} con cronograma
            anunciado y compensación que mitigue el impacto sobre los hogares vulnerables.
      \item \alert{Vigilar / Sin presión inmediata:} mantener monitoreo. Un nuevo choque de
            precios puede desplazar a estos países hacia el cuadrante de reforma urgente,
            dado el carácter \textbf{pro-cíclico} del subsidio explícito.
    \end{itemize}
  \end{block}
  \vspace{-0.2em}
  {\footnotesize \textbf{Secuencia transversal:} (1) congelar la indexación del subsidio al precio
  internacional; (2) liberar el precio de forma gradual y anunciada; (3) reasignar el gasto a
  \textbf{transferencias focalizadas} a los hogares vulnerables.}

  \note{%
    Este es el slide que convierte el coeficiente y la matriz en acción concreta, que es lo que el
    reto pide. La idea central: la recomendación se diferencia por cuadrante. En reforma urgente no
    hay margen (los déficit ya son fuertes, Bolivia y Venezuela), hay que actuar ya con compensación
    inmediata. En reforma gradual hay colchón de deuda baja, así que se puede secuenciar con
    cronograma. En los otros dos, vigilar, porque el subsidio es pro-cíclico y un nuevo choque los
    mueve al cuadrante malo. Y doy una secuencia operativa de tres pasos: congelar indexación,
    liberar precio gradual, redirigir el ahorro a transferencias focalizadas. La focalización (no
    universal) viene de la regresividad del slide anterior.}
\end{frame}

% --- Slide 12 (existente): Conclusión y recomendación ---
\begin{frame}{Conclusión y recomendación de política}
  \footnotesize
  \begin{itemize}\setlength{\itemsep}{2pt}
    \item \textbf{Hallazgo:} el choque de 2022 incrementó \textbf{diferencialmente} el subsidio
          explícito de los exportadores netos en \key{+1.80 pp del PIB}: efecto robusto en signo
          y magnitud (estable bajo leave-one-out), significativo a $\alpha = 0.10$
          ($p \approx 0.06$), concentrado en el componente sensible al precio y \textbf{transitorio}
          (se revierte en 2023, cuando el Brent cae).
    \item \textbf{Alcance:} la vulnerabilidad fiscal del subsidio \textbf{precede al choque} y
          alcanza también a importadores (Surinam, Argentina). El choque \textbf{agravó} una
          exposición preexistente, no la creó.
  \end{itemize}

  \begin{block}{Recomendación: reforma diferenciada por espacio fiscal (no uniforme)}
    \begin{itemize}\setlength{\itemsep}{1pt}
      \item Subsidio alto y deuda alta: \alert{reforma prioritaria} con compensación
            focalizada inmediata.
      \item Subsidio alto y margen fiscal disponible: \textbf{transición gradual} con
            cronograma anunciado y transferencias focalizadas a hogares vulnerables.
    \end{itemize}
    \vspace{0.1em}
    {\footnotesize El subsidio universal al combustible es \textbf{regresivo} — el quintil
    más rico capta más del 40\% de los beneficios — por lo que sustituirlo por transferencias
    focalizadas mejora simultáneamente la equidad y el balance fiscal.}
  \end{block}
  \vspace{-0.3em}
  \begin{center}\color{Carbon}\small
  Un subsidio que crece con cada choque de precios constituye un \textbf{pasivo fiscal
  pro-cíclico}: reformarlo libera espacio fiscal y reduce la exposición a la volatilidad
  del precio del petróleo.
  \end{center}

  \note{%
    Cierro uniendo evidencia y política. El hallazgo en una frase: el choque afecta con más fuerza a
    los exportadores, pero de forma transitoria. El matiz evita la lectura simplista: no es solo un
    problema de exportadores; el subsidio elevado ya constituía un riesgo fiscal en varios
    importadores. La recomendación no es ``eliminar subsidios'' en abstracto, sino diferenciar por
    espacio fiscal. La idea central: el subsidio es pro-cíclico, crece justo cuando el petróleo sube,
    que es cuando más cuesta sostenerlo; reformarlo aumenta la resiliencia fiscal. Ese es el mensaje
    de política central.}
\end{frame}

% ============================================================
% REFERENCIAS
% ============================================================
\begin{frame}{Referencias}
  \scriptsize
  \begin{itemize}\setlength{\itemsep}{1pt}
    \item Angrist, J. D., \& Pischke, J.-S. (2009). \emph{Mostly Harmless Econometrics}.
          Princeton University Press. \textbf{[DiD 2$\times$2 canónico]}
    \item Black, S., Liu, A. A., Parry, I., \& Vernon, N. (2023). \emph{IMF Fossil Fuel Subsidies
          Data: 2023 Update}. IMF Working Paper No.\ 2023/169.
    \item Coady, D., Flamini, V., \& Sears, L. (2015). \emph{The Unequal Benefits of Fuel Subsidies
          Revisited}. IMF Working Paper WP/15/250.
    \item Roth, J., Sant'Anna, P., Bilinski, A., \& Poet, J. (2023). What's trending in
          difference-in-differences? \emph{Journal of Econometrics}.
    \item International Monetary Fund. \emph{Fossil Fuel Subsidies Database} (2023 update).
          \textbf{[Fuente de datos principal]}
    \item International Monetary Fund. \emph{World Economic Outlook Database}
          (IMF DataMapper API). Indicadores: \texttt{GGXWDG\_NGDP}, \texttt{GGXCNL\_NGDP},
          \texttt{GGR\_G01\_GDP\_PT} (\% del PIB). \textbf{[Variables fiscales]}
    \item U.S. Energy Information Administration. \emph{Europe Brent Spot Price}.
          \textbf{[Precio del petróleo]}
  \end{itemize}

  \note{%
    No leo la lista; es respaldo. Las anclas son tres: el contexto del choque y la regresividad
    (FMI: Black et al.\ 2023 y Coady et al.\ 2015), el marco DiD y sus supuestos (Roth et al.\ 2023),
    y las fuentes de datos (IMF FFS Database y EIA). Mantengo la bibliografía corta a propósito: es
    una prueba técnica, no una tesis.}
\end{frame}

% ============================================================
% (El antiguo "Material de respaldo" se eliminó: sus 4 piezas
%  — tab1, tab2, fig3, tab5 — se integraron al cuerpo principal
%  con lectura guiada en sus bloques respectivos.)
% ============================================================

% ============================================================
% RESPALDO — PREGUNTAS ANTICIPADAS (Q&A de defensa)
\end{document}
